\documentclass[11pt]{article}
\usepackage[a4paper,margin=26mm]{geometry}
\usepackage[T1]{fontenc}
\usepackage{lmodern}
\usepackage{amsmath,amssymb,amsthm}
\usepackage{microtype}
\usepackage[hidelinks]{hyperref}
\newtheorem{theorem}{Theorem}
\newtheorem{lemma}{Lemma}
\newtheorem{proposition}{Proposition}
\newcommand{\Q}{\mathbb{Q}}
\newcommand{\Z}{\mathbb{Z}}
\newcommand{\N}{\mathbb{N}}
\newcommand{\val}{\operatorname{val}}
\title{A continuous valuation with rational value group\\
\large Disproof of Conjecture 00000005754}
\author{}
\date{}
\begin{document}
\maketitle

\begin{abstract}
The Hahn-series field with rational coefficients and rational exponents carries a continuous
valuation whose intrinsic value group is the additive group $\Q$. This group is not isomorphic
to $\Z^I$ for any index set $I$, even without requiring preservation of order. Consequently no
choice of rank can make the conjecture's asserted integer-power description hold. The Lean
formalization constructs the valuation, proves continuity for two explicit codomain topologies,
identifies the group of values actually attained, and proves the obstruction for every index type.
\end{abstract}

\section{Claim being disproved and conventions}

The final clause of the supplied English statement says that ``the value group of a continuous
valuation is a power of Z, with exponent the rank.'' The Chinese statement makes the same
assertion. This is a universal assertion about continuous valuations, and it is a conjunct of the
conjecture. One continuous valuation whose value group is no integer power therefore suffices
to disprove the conjecture as written.

The statement gives neither a topology nor a discreteness hypothesis. We use the standard
valuation topology on the domain field. We prove continuity both for the usual order topology
on the extended value group and for the finer topology that isolates finite additive values and
has tails as neighborhoods of infinity. Thus the counterexample does not rely on assigning the
domain an arbitrary discrete topology. A field is, in particular, a ring of the kind mentioned in
the statement.

We use additive valuation notation. An additive valuation on a field $K$ is a function
\[
 w:K\longrightarrow \Gamma\cup\{\infty\}
\]
into a linearly ordered abelian group with a largest element adjoined, satisfying
\[
 w(0)=\infty,\qquad w(1)=0,\qquad
 w(fg)=w(f)+w(g),\qquad w(f+g)\geq\min\{w(f),w(g)\}.
\]
Here $\infty+\gamma=\infty$. Its intrinsic additive value group is the image of $K^\times$
under $w$, with the induced addition and order. In particular, an unnecessarily large target
group does not enlarge this intrinsic group.

No interpretation of the unspecified conversion formula or of the asserted rank jumps is
needed. We will prove the stronger statement that the intrinsic group is not $\Z^I$ for
\emph{any} index set $I$. In particular, it cannot be $\Z^{r}$ when $r$ is whatever rank is
intended in the statement. When $r\in\N$, the corresponding Lean type is
\texttt{Fin r $\to$ $\Z$}.

\section{The field and its valuation}

Let $K$ be the Hahn-series field
\[
 K=\Q((t^{\Q})).
\]
Its elements are formal sums
\[
 f=\sum_{q\in\Q}a_qt^q,\qquad a_q\in\Q,
\]
whose support $\{q:a_q\ne0\}$ is well ordered. Addition is coefficientwise and multiplication
is convolution. The standard Hahn-series construction makes this a field. In particular,
monomials $t^q$ are nonzero, satisfy $t^qt^r=t^{q+r}$, and have inverse $t^{-q}$.
The field structure, including the general inversion construction, is provided and proved by
Mathlib's \texttt{HahnSeries.instField} in
\texttt{Mathlib/RingTheory/HahnSeries/Summable.lean}.

Define
\[
 w(0)=\infty,\qquad
 w(f)=\min\operatorname{supp}(f)\quad(f\ne0).
\]
The minimum exists by well ordering. In a product of two nonzero series, the coefficient at
$w(f)+w(g)$ is the product of the two leading coefficients. It is nonzero, and no smaller
exponent occurs. Therefore $w(fg)=w(f)+w(g)$. Addition cannot create an exponent below both
leading exponents, so $w(f+g)\geq\min\{w(f),w(g)\}$. Also $w(1)=0$.
This is the additive valuation \texttt{HahnSeries.addVal $\Q$ $\Q$} in the formalization.

For every $q\in\Q$,
\begin{equation}\label{eq:monomial}
 t^q\ne0,\qquad w(t^q)=q.
\end{equation}
Thus every rational value is actually attained on a nonzero field element.

\section{Topology and continuity}

The valuation topology on $K$ has the following neighborhood basis at $a\in K$:
\[
 B_\gamma(a)=\{f\in K:w(f-a)>\gamma\},\qquad \gamma\in\Q.
\]
This is the canonical Hausdorff topological-field structure associated with the valuation.
The Lean instance is built with \texttt{Valued.mk'}; its topological division-ring and
Hausdorff instances are also checked in \texttt{Inspect.lean}.

First give $\Q\cup\{\infty\}$ the topology that isolates every rational value and has
$\{q:q>\gamma\}\cup\{\infty\}$ as a neighborhood basis at $\infty$.

\begin{lemma}
The valuation $w$ is continuous for this topology on its codomain.
\end{lemma}
\begin{proof}
At zero, the inverse image of the indicated tail is exactly $B_\gamma(0)$.
For $f\ne0$, put $r=w(f)$. If $g\in B_r(f)$, then $w(g-f)>r$.
The valuation inequality gives $w(g)\geq r$. If $w(g)>r$, applying the same inequality to
$f=g-(g-f)$ would give $w(f)>r$, a contradiction. Therefore $w(g)=r$ throughout
$B_r(f)$. The valuation is locally constant at every nonzero point, and hence continuous there.
\end{proof}

The ordinary order topology on $\Q\cup\{\infty\}$ is coarser than this topology.
Indeed, finite points are isolated in the finer topology, and order neighborhoods of infinity
contain tails. Hence $w$ is continuous also for the ordinary order topology.

Mathlib's multiplicative notation reverses the additive order: in Lean the extended target is
\[
 \texttt{ExtendedValue = Multiplicative (WithTop $\Q$)$^{\mathrm{od}}$}.
\]
Its multiplicative zero denotes additive infinity, and multiplication denotes addition of
exponents. The definition \texttt{valuation} is exactly
\texttt{additiveValuation.toValuation}; it does not alter the underlying values.
The finer codomain topology is \texttt{WithZeroTopology}. The formal theorem
\texttt{valuation\_continuous} uses this topology, while
\texttt{valuation\_continuous\_orderTopology} explicitly uses
\texttt{Preorder.topology ExtendedValue}. Reversing a linear order interchanges the two kinds
of open rays, and therefore leaves the ordinary order topology unchanged. These theorems thus
establish the two continuity statements above in the multiplicative encoding.

\section{The intrinsic value group}

The formalization defines an additive-group homomorphism
\[
 \texttt{valueHom}:\operatorname{Additive}(K^\times)\longrightarrow\Q,
 \qquad u\longmapsto \min\operatorname{supp}(u).
\]
The additive wrapper means that multiplication of units is regarded as addition. The identity
$w(uv)=w(u)+w(v)$ is exactly the homomorphism law. The intrinsic group is defined as
\[
 G=\operatorname{range}(\texttt{valueHom})\leq(\Q,+),
\]
called \texttt{valueGroup} in Lean.

The theorem \texttt{valueHom\_agrees} identifies this homomorphism with the finite part
of the valuation on units. The theorem \texttt{mem\_valueGroup\_iff} gives the following
connection to nonzero field elements:
\[
 q\in G\quad\Longleftrightarrow\quad
 \exists f\in K,\ f\ne0\ \text{and}\ w(f)=q.
\]
Equation~\eqref{eq:monomial} proves $G=\Q$, formally
\texttt{valueGroup\_eq\_top}. The resulting additive isomorphism is
\texttt{valueGroupEquiv}. Thus the subsequent group obstruction concerns the actual image
group of the constructed valuation.

\section{No integer power is possible}

For an index set $I$, write $\Z^I$ for the group of functions from $I$ to $\Z$ with pointwise
addition.

\begin{proposition}\label{prop:obstruction}
For every index set $I$, there is no additive-group isomorphism $\Q\cong\Z^I$.
\end{proposition}
\begin{proof}
Suppose $e:\Q\to\Z^I$ were such an isomorphism.
If $I$ is empty, $\Z^I$ is the trivial group, whereas $0\ne1$ in $\Q$, contradicting
injectivity.

If $I$ is nonempty, fix $i\in I$ and let $\mathbf{1}$ denote the function that is identically
one. Set $x=e^{-1}(\mathbf{1})$ and $y=x/2$. Since $y+y=x$, additivity gives
\[
 e(y)+e(y)=\mathbf{1}.
\]
Evaluating at $i$ yields $2e(y)(i)=1$ in $\Z$, which is impossible. This contradiction
proves the assertion.
\end{proof}

The formal theorem \texttt{rationals\_not\_integer\_power} proves
Proposition~\ref{prop:obstruction} for an arbitrary Lean index type. Composing a proposed
isomorphism $G\cong\Z^I$ with \texttt{valueGroupEquiv} then proves
\texttt{valueGroup\_not\_integer\_power}. This rules out even abstract additive isomorphisms,
so it also rules out ordered-group isomorphisms.

\begin{theorem}
There is a continuous valuation on a Hausdorff topological field whose intrinsic value group
is not a power of $\Z$. Consequently Conjecture 00000005754 is false as stated.
\end{theorem}
\begin{proof}
Take the field and valuation constructed above. Continuity holds for both explicit codomain
topologies. Its intrinsic value group is $G\cong\Q$, and
Proposition~\ref{prop:obstruction} rules out $G\cong\Z^I$ for every $I$.
In particular, for no natural number $r$ does $G\cong\Z^r$ hold. The final universal clause
of the conjecture therefore fails for this example, independently of the other clauses.
\end{proof}

\section{Formalization and verification scope}

The mathematical source is \texttt{Counterexample.lean}, using Lean 4.19.0 and Mathlib
v4.19.0 pinned to commit
\begin{center}
\texttt{c44e0c8ee63ca166450922a373c7409c5d26b00b}.
\end{center}
The theorem \texttt{continuous\_valuation\_counterexample} combines continuity for both
topologies, attainment of every rational value by a nonzero series, the obstruction for every
index type, and the explicit negation
\[
 \neg\exists r\in\N,\ G\cong(\operatorname{Fin}(r)\to\Z).
\]
The exact definition of $G$ and its relation to $w$ are proved in the preceding named theorems.
No unproved rank assignment is assumed.

An independent fresh project build with \texttt{lake build}, direct replay of the mathematical
source with warnings treated as errors, and a compiled-declaration audit have passed.
The audit in \texttt{lean/Audit.lean} enumerates every declaration owned by the mathematical
module, including compiler-generated declarations: 24 in total. It checks each declaration's
transitive axiom dependencies, type, and direct constants. Only the standard foundational
axioms \texttt{propext}, \texttt{Classical.choice}, and \texttt{Quot.sound} occur.
\texttt{lean/Inspect.lean} displays the exact theorem types, including topology arguments,
and checks the field, Hausdorff, and topological division-ring instances.
There are no custom axioms, proof placeholders, or untrusted native decision proofs.
No auxiliary numerical or symbolic computation is required by the argument.

The supplied \texttt{verify.py} rebuilds the authored project in a new directory, replays its
sources with warnings treated as errors, compares the complete compiled inventory with the
frozen reviewed inventory, and checks the nine dependency revisions and source cleanliness
before and after execution. It records actual commands, exit codes, and output. The companion
\texttt{test\_verify.py} contains verifier controls, including compiled examples that must be
rejected for an added axiom, an unsafe definition, or an admitted proof. Reproduction commands
and the validation records are listed in \texttt{README.md} and \texttt{VERIFICATION.md}.
These are local checks and an independent local semantic review; maintainer acceptance is a
separate repository decision.

The supplied original statement was verified against its SHA-256 identity:
\begin{center}\small
\texttt{e9168d9270246842f8d39c8e5f8dae0ff58c}\newline
\texttt{4723f9957975245e4ab785b720f6}.
\end{center}
The access-provenance and verification logs record the isolated inputs, consulted standard
library sources, dependency identities, actual build commands, and their results.

\end{document}
